\chapter{绪论}

% 是论文正文的开端，应包括毕业论文选题的背景、目的和意义；
% 对国内外研究现状和相关领域中已有的研究成果的简要评述；
% 介绍本项研究工作研究设想、研究方法或实验设计、理论依据或实验基础；涉及范围和预期结果等。
% 要求言简意赅，注意不要与摘要雷同或成为摘要的注解。
\label{cha:introduction}
\section{选题背景与意义}
\label{sec:background}
人机运动重定向是指将人类动作序列从人的关节运动空间映射到机器人运动空间的过程，
旨在解决两者的物理结构、运动学约束以及自由度配置之间的差异。
该技术可根据应用需求分为实时重定向与离线重定向：实时重定向是机器人遥操作技术的核心，
广泛应用于远程手术等需要精密操作的场景；而离线重定向则在机器人数据采集领域扮演着愈发重要的角色。
随着具身智能的飞速发展，高质量、规模化的机器人动作数据已成为驱动大模型训练的基础，因此机器人数据采集工作
也变得变得重要。

而在对于机器人的数据采集来说，灵巧手的数据精度要求更高。
同时，灵巧手自由度远比其他部位高，使得重定向算法在处理空间映射时面临极大的复杂性。
在机器人系统中，灵巧手作为与环境交互的主要末端执行器，其控制精度直接影响任务的完成质量。
另一方面，由于不同厂商生产的灵巧手在机械构造和自由度上存在显著异构性，
开发一种具有平台兼容性且能适配多种灵巧手模型的通用算法具有重要的工程价值。

目前，现有的优化方法或基于学习的方法虽然在速度或特定任务的精度上取得了进展，
但大多仅关注单帧数据的映射，忽略了连续输入帧之间的时序相关性。
这种局限性导致生成的机器人动作序列在物理执行时缺乏平滑性，极易受到视觉捕捉噪声的干扰而产生抖动。
在灵巧手的运动重定向之中，应该要考虑时序信息，以提升重定向的精度。
因此，研究并开发一种能够深度提取输入序列时空特征、具备高精度、强实时性且兼顾泛化性能的运动重定向框架，
对于推动具身智能的数据获取效率及提升灵巧遥操作系统的稳定性具有重要的理论意义与应用前景
\cite{oneillOpenXembodimentRobotic2024}\cite{brohanRt1RoboticsTransformer2022}。

\section{国内外研究现状和相关工作}
\label{sec:related_work}
\subsection{基于优化的运动重定向方法}
早期的运动重定向研究主要依赖于数值优化算法。
这类方法通过建立源域与目标域之间的运动学映射方程，并在逆运动学框架下施加关节限位、自碰撞检测及平滑度约束，
寻找每个时刻的最优关节解。例如，Qin等人提出的AnyTeleop框架利用统一的运动学映射实现了较好的实时性与平台泛化能力。
然而，基于优化的方法通常依赖于简化的数学模型，在处理高度异构的非线性映射时，容易陷入局部最优解，
导致在复杂抓取任务中的重定向精度受限。

\subsection{基于学习的运动重定向方法}
随着深度学习技术的演进，越来越多的研究者转向数据驱动的方法。通过构建深层神经网络，
模型能够直接学习从人手姿态到机器人关节空间的复杂非线性映射。这类方法通过增加模型参数量，
显著提升了重定向的精度与对噪声的鲁棒性。例如，Figuera等人通过重新定义数据配对方式，强化了模型对精细动作的捕捉
\cite{figueraRedefiningDataPairing2024}。
然而，纯学习方法通常面临泛化性挑战，即针对特定机器人型号训练的模型难以直接部署到其他异构平台上。
此外，随着模型规模的增大，其推理延迟往往会增加，给实时遥操作任务带来挑战。

\subsection{优化与学习结合的混合方法}
为了兼顾重定向的速度、精度与通用性，近年来的研究开始探索将学习结合优化的混合框架。
这类方法通常利用神经网络快速预测初始姿态，再结合优化层进行运动学约束校准，
从而在保证精度的同时提升跨平台的适配能力\cite{zhangKinematicMotionRetargeting2022}。
尽管如此，上述方法大多将重定向任务视作单帧映射问题，
在实际运行过程中未能有效提取连续输入帧之间的时空相关信息。由于缺乏帧间依赖关系的建模，
生成的机器人动作序列在面对视觉捕捉抖动或环境噪声时，往往缺乏必要的平滑度，导致物理执行器的机械磨损加剧且操作体验不佳。

\subsection{Transformer在时序特征提取中的应用}
Transformer架构凭借其强大的自注意力机制和并行计算能力，
在自然语言处理及计算机视觉领域取得了突破性进展\cite{vaswaniAttentionAllYou2017}。在智能机器人方面，Transformer 已开始应用于人体骨骼序列建模，
用于捕捉复杂的运动轨迹规律。相比于传统的循环神经网络（RNN），Transformer 能够更有效地处理长距离依赖关系。
另外也有研究使用Transformer架构来将二维人体姿态序列升维至三维人体序列，展示了其对三维姿态序列的建模能力。
受此启发，在本研究中，我们引入时空 Transformer 架构来处理连续的人手三维运动序列，
旨在通过建模帧间的时间依赖性与帧内的空间关联性，从根本上解决单帧重定向带来的动作不稳定问题,提高重定向轨迹的平滑度。
% \footnote{大家要养成添加脚注的好习惯}
% \subsection{基于统计的颜色映射方法}
% 两级标题之间要有过渡性文字。可以通过一段话引出下面的文字或者对本章内容概括。
% 论文中凡非正式参考文献以外的资料，应以脚注的方式注明。

\section{本文的研究内容与主要工作}
本文旨在针对灵巧机器人动作重定向中存在的时空特征提取不足、动作不平滑等问题，
提出并实现一种基于时空 Transformer 架构的通用重定向框架 TransRetarget。具体研究内容与主要工作如下：

\subsection{基于时空 Transformer 的重定向模型架构设计}
本文设计并实现了一种全新的多帧输入重定向网络。通过引入 Transformer 编码器，
利用自注意力机制同时对人手动作序列进行空间维度（帧内关节相关性）和时间维度（帧间时序依赖性）的特征建模。
该架构改变了传统方法逐帧推理的模式，能够有效识别并抑制输入信号中的瞬时噪声与抖动，从而生成高平滑度的机器人动作指令。

\subsection{面向自监督学习的多准则复合损失函数构建}
为了实现高精度的动作对齐，本文设计了一套综合运动学特征约束的复合损失函数系统。
该系统整合了手指侧向移动损失、手指弯曲度一致性损失以及关键指尖间距约束。通过在特征空间内最小化人手向量与机器人手向量的差异，
强制模型学习到具有物理语义的一致性映射，显著提升了生成动作的自然度与抓取准确性。

\subsection{跨平台泛化性与实时控制系统实现}
本文构建了一个具备高通用性的软件框架，能够适配多种异构灵巧手模型。
通过引入归一化参数，解决了不同硬件间自由度配置与物理比例不一致的问题。
同时，针对实时遥操作场景优化了模型推理逻辑，
确保了算法在保持高精度的前提下，具备极低的计算延迟，满足实时交互控制的需求。

\subsection{实验验证与性能评估}
本文在一个公开数据集和一个使用实验室器材自行采集的VisionPro动捕数据集上进行了充分的实验验证。通过与现有基于优化、
基于学习以及混合策略的基准算法进行对比，从重定向误差、轨迹平滑度、计算效率及跨平台成功率等多个维度评估了 TransRetarget 的优越性。
此外，本文还开展了消融实验，验证了各项损失函数对系统最终性能的贡献。

\section{本文的论文结构与章节安排}
\label{sec:arrangement}
本文共分为五章，各章节内容安排如下：

第一章为绪论，介绍了研究背景、意义、国内外研究现状以及本文的主要工作和章节安排。

第二章为相关工作，详细回顾了运动重定向领域的经典方法与最新进展，并分析了现有方法的优缺点。同时介绍了Transformer
架构的基本原理与实现方法，以及其在人体骨骼序列中的应用。

第三章为方法设计，在本章我们系统阐述了TransRetarget这个方法，介绍了基于时空 Transformer 的重定向模型架构、损失函数设计。

第四章为实验与结果，展示了在多个数据集上的实验设置、评估指标以及与基准方法的对比结果，并通过消融实验分析了各模块的贡献。

第五章为重定向结果的交互强化，探讨了如何在得到重定向结果后，在实时交互的同时进行合理修正，
以强化交互效果。

第六章为总结与展望，总结了本文的研究工作与主要贡献，并对未来的研究方向进行了展望。

